\section{Diagnostics and tests}

Evaluar si un método de SBI está funcionando bien es clave, porque a diferencia de métodos más clásicos como MCMC, las redes neuronales pueden fallar de forma silenciosa sin que sea obvio. No hay algo como un diagnóstico de convergencia estándar que sea seguro, así que hay que hacer tests específicos para ver si el posterior que aprendimos realmente tiene sentido. La idea general de estos diagnósticos es chequear consistencia estadística usando el simulador, probando el método en muchos datos simulados donde sabemos la verdad.

\subsection{Posterior predictive checks (PPC)}

Uno de los chequeos más básicos es el \textit{posterior predictive check} (PPC). La idea es simple: se toman muestras de parámetros desde el posterior estimado $p(\Theta \mid x_0)$, luego simulas nuevos datos usando el modelo $p(x \mid \Theta)$, y comparas esos datos simulados con los datos observados. Matemáticamente, esto corresponde a mirar la distribución predictiva:
\[
p(x \mid x_0) = \int p(x \mid \Theta)\, p(\Theta \mid x_0)\, d\Theta.
\]
Si los datos simulados se parecen a los observados, es una señal de que el modelo está capturando bien el proceso generativo. Eso sí, pasar el posterior predictive check no garantiza que el posterior sea bueno. Por ejemplo, un posterior muy poco informativo podría generar datos variados que “cubren” los datos observados y aun así no decir nada útil sobre los parámetros. Entonces, este chequeo es necesario, pero no suficiente: sirve como filtro básico, pero hay que complementarlo con otros tests más exigentes.

\subsection{Coverage diagnostics}

Una forma más exigente de validar los posteriors es a través de \textit{coverage diagnostics}. La idea acá es ver si los intervalos de credibilidad que entrega el modelo realmente contienen el valor verdadero con la frecuencia correcta, en promedio. Para esto se construye un conjunto de calibración generando pares $(\Theta_i, x_i)$ desde la distribución conjunta $p(\Theta, x)$, de modo que conocemos el parámetro verdadero que generó cada observación.

Primero, uno de los enfoques mas interpretables es el de \textit{expected coverage} \cite{hermans2022}. En este caso se evalúa el posterior completo usando regiones de alta densidad (HPD). La idea es ver si, por ejemplo, una región de credibilidad del $90\%$ realmente contiene el valor verdadero el $90\%$ de las veces. Esto se suele visualizar con gráficos tipo P-P plot o Z-Z plot: si el modelo está bien calibrado, aparece una diagonal. Si el posterior es muy conservador, la curva queda arriba; si es muy confiado, queda abajo. Este método es potente porque evalúa la distribución conjunta y puede detectar problemas en correlaciones entre parámetros. Luego está TARP (\textit{Tests of Accuracy with Random Points}) \cite{lemos2023}, que usa una idea distinta: construye regiones alrededor de puntos de referencia aleatorios en el espacio de parámetros y compara distancias entre muestras del posterior y el valor verdadero. En términos simples, verifica si el parámetro real cae lo suficientemente cerca de muestras típicas del posterior. Este enfoque permite detectar ciertos sesgos que métodos basados en regiones HPD pueden pasar por alto, y además es útil cuando no se puede evaluar explícitamente la densidad del modelo. Finalmente, está \textit{Simulation-Based Calibration} (SBC) \cite{talts2020}, que es más simple y se basa en un procedimiento de ranking. Para cada ejemplo del conjunto de calibración, se toman varias muestras del posterior y se compara el valor verdadero con ellas, viendo en qué posición queda. Si el posterior está bien calibrado, estos rankings deberían distribuirse uniformemente.

Algo importante es que todos estos tests son computacionalmente caros, porque requieren hacer inferencia muchas veces. Métodos como NPE son más rápidos para esto porque generan muestras directamente, mientras que NLE o NRE suelen necesitar MCMC en cada caso, lo que los hace más lentos. Además, en métodos secuenciales como SNPE, estos tests se complican más porque los posteriors aprendidos dependen de la proposal y no son válidos globalmente. Finalmente, pasar estos tests es una condición necesaria, pero no suficiente. Por ejemplo, un posterior muy ancho puede estar perfectamente calibrado en términos de coverage, pero no ser útil en la práctica. Así que, como siempre, hay que interpretar estos diagnósticos con cuidado.